% LSPV-5 FACTOREDGATE1, 2026-11-18 07:06 PST. Substantive standalone GitHub manuscript.
\documentclass[11pt]{article}
\usepackage[margin=0.85in]{geometry}
\usepackage{amsmath,amssymb,amsthm,hyperref}
\newtheorem{theorem}{Theorem}
\newtheorem{proposition}[theorem]{Proposition}
\title{FACTOREDGATE1: Source-Discovered Conjugacy and the Semantic Verification Gate}
\author{LSPV-5 Source Certificates}
\date{18 November 2026}
\begin{document}\maketitle
\noindent\textbf{Status.} Restricted construction and coNP reduction PROVED/EXACT subject to independent review; finite tests COMPUTATIONAL; universal two-output PFCERT CONDITIONAL/OPEN, nonalgebrization OPEN, and $P\ne NP$ OPEN. Full six-page manuscript, bibliography, implementation and immutable raw checkpoints are in the locally preserved FACTOREDGATE1 research bundle.

\section{IDEAS PROPOSED BEFORE LITERATURE}
Original four CANDIDATE proposals were SHA-256 frozen \emph{before} targeted literature search in \texttt{tests/ideas\_before\_literature.txt}, hash \texttt{6d160e5cd81eca46b8f8be0d632bb3ac6ffadee83bd81bff4d97d70a8b05bd8c}.
CANDIDATE A proposed extracting a hidden $F=H^{-1}A_bH$ directly from an involutive source-gate word using only local cancellations, with invalid inverse order and forged output as falsifiers. Missing premise was sound verification without general semantic circuit equivalence. CANDIDATE B proposed a polynomial-size conjugated source with exponentially many ANF monomials; missing premise was an exact support count; full truth-table M\"obius transforms and orbit checks were falsifiers. CANDIDATE C proposed coNP-hardness of semantic identity checking even at $H=A_b=I$, falsified if the UNSAT reduction failed. CANDIDATE D proposed including singular affine cores, falsified by singular-core endpoint disagreement. Original proposals remain separate from the following post-search revisions.

\section{LITERATURE CHECK AND REVISIONS}
Bennett (IBM JRD 1973, DOI 10.1147/rd.176.0525), Toffoli (MIT/LCS/TM-151, 1980), and Miller--Maslov--Dueck (DAC 2003, DOI 10.1145/775832.775915) support using gate-level reversible composition but do not entail fast simulation for arbitrary reversible machines. Finite-field Koopman observability work (Anantharaman--Sule, arXiv:2404.02336) provides the existing linear-observable framework; KOOPMON1's unbounded monomial expansion motivates B. General circuit tautology/equivalence checking supplies the coNP gate, and motivates restricting A to sound syntactic rewriting. Hirahara, ECCC TR20-050 (2020), Theorem 3.9 uses \emph{plain} time-bounded $K^t$, not Levin $Kt$. The original proof at pp. 29--30 uses an $O(\log n)$-suffix search. Baker--Gill--Solovay (SIAM J. Comput. 4(4), 1975, DOI 10.1137/0204037) use deterministic and nondeterministic \emph{oracle query machines}; the accessible publisher article abstract verifies the model and opposite-answer oracle existence. Aaronson--Wigderson (TOCT 2009) impose a separate algebrization check. Original Drive NONALG-0003 and PVG-4-0001 were read. Revisions: A becomes a sound but incomplete syntactic recognizer; B, C and D survive in restricted scope. No literature absence is construed as proof of novelty.

\section{CONSTRUCTION OR PROOF ATTEMPT}
\paragraph{Authentic source and discoverable normal form.}
A canonical source of length $N$ lists $n$ Boolean state coordinates, explicit initial state $x_0$, binary horizon $T$ with $L$ bits, one output coordinate, a source \emph{prefix} word $U$ of controlled-XOR gates, an explicit affine core $A_b(x)=Ax+b$ with all $n$ rows, and a source \emph{suffix} word $V$. A gate changes its target $i$ by the AND of listed literals on distinct coordinates other than $i$; hence each gate is an involution. Every index and control and the whole matrix are part of the source; $n,L\le N$. A deterministic stack cancels adjacent equal gates to give $\mathrm{red}(U),\mathrm{red}(V)$. It admits the source only if
\[
\mathrm{red}(V)=\mathrm{reverse}(\mathrm{red}(U)).
\]
Then $H$ executes $\mathrm{red}(U)$ and the \emph{all-state} source identity $F(x)=H^{-1}A_bH(x)$ follows by locally justified cancellation. A valid source outside the discoverable grammar yields \texttt{NOT\_CERTIFIED}, not UNSAT and not a lower bound. This is not supplied quotient access or semantic PIT.

\begin{theorem}[Restricted polynomial endpoint certificate]
For each valid authenticated source satisfying the displayed syntactic rule, a deterministic verifier recomputes the extracted $H$, the augmented matrix $M=\left(\begin{smallmatrix} A&b\\0&1\end{smallmatrix}\right)$ and the sequence $P_0=M,\ P_{j+1}=P_j^2$ through the binary horizon. It verifies $z_0=(H(x_0),1)$ and $z_T=(\prod_{j:T_j=1}P_j)z_0$, then returns $H^{-1}(\mathrm{head}_n(z_T))$ at the selected output coordinate. Certificate length is $O(N+(n+1)^2\max(1,L))$ bits and the complete parser, discovery, certificate constructor and verifier cost $\operatorname{poly}(N)$ bit time and simultaneous peak space. Initiality, all-state consecution, terminal coverage and source authentication are independently checked. Every accepted certificate on the same source yields \emph{exactly one global output bit} across all accepted proofs.
\end{theorem}
\begin{proof}
Every gate cancels with its identical neighbor and is its own inverse, giving $F=H^{-1}A_bH$ as functions on all states, not sampled trajectories. Consequently the transformed trajectory obeys $(z_{t+1},1)=M(z_t,1)$ for every state. Induction and exact binary exponentiation yield the endpoint. Every power square, physical output and source identity is checked against the original canonical source, so no witness choice changes the output. The verifier stores at most $O((n+1)^2L+N)$ source/certificate bits and performs $O((n+1)^3L)+\operatorname{poly}(N)$ elementary bit operations. $A$ need not be invertible.
\end{proof}

\begin{proposition}[Exact exponential symbolic expansion is not necessary work]
For $k\ge2$, use $x=(a,z_1,\dots,z_k)$, one reversible source gate $H(a,z)=(a+\prod_{i=1}^k(1+z_i),z)$ and $A_b(a,z)=(a,z+\mathbf1)$. The explicit full-matrix encoding has $O(k^2+k\log k)$ bits. The first coordinate of $F=H^{-1}A_bH$ is
\[
F_0(a,z)=a+\prod_i(1+z_i)+\prod_i z_i
=a+\sum_{S\subsetneq[k]} \prod_{i\in S}z_i.
\]
Its unique Boolean ANF has exactly $2^k$ monomials and the substitution-closed monomial family seeded by this output has exactly $2^k$ features. Nevertheless $F^2=I$ and any $F^T(x_0)$ has an $O(\operatorname{poly}(k+\log T))$ factored evaluator.
\end{proposition}
\begin{proof}
In characteristic two the full-degree monomial cancels, all proper-subset monomials remain. Substituting $z_i\mapsto z_i+1$ sends proper subsets to their subsets, so this monomial family is closed. Both $H^2$ and $A_b^2$ are the identity, so $F^2=I$. The certified source algorithm avoids expansion.
\end{proof}

\begin{theorem}[General semantic verification is coNP-complete]
Given an arbitrary Boolean circuit $F$, an explicitly evaluable reversible gate-word $H$, and an affine $A_b$, verifying $\forall x:F(x)=H^{-1}A_bH(x)$ is coNP-complete, already for $H=A_b=I$.
\end{theorem}
\begin{proof}
A nonidentity input is an NP witness against equality. For CNF $\varphi(u)$ set $F_\varphi(u,b)=(u,b\oplus\varphi(u))$. Then $F_\varphi=I$ iff $\varphi$ is UNSAT. The mapping is polynomial. This verifies an exact \emph{semantic checking} gate, not a time lower bound.
\end{proof}
Syntactic completeness fails even for a two-bit word NOT(0),NOT(1),NOT(0),NOT(1), which is semantically the identity by commuting flips but does not reduce through adjacent-pair cancellation. Thus soundness is proved; discovery of arbitrary equivalent structures remains OPEN.

\paragraph{Matched-oracle diagnostic and bounded-output semantics.}
For any \emph{one fixed} oracle $O$, choose a universal oracle-query machine $U^O$ with efficient fixed-machine simulation. Define $K_{U}^{t,O}(x)$ as the shortest program outputting $x$ within $t(|x|)$ oracle-machine steps, and let $R_{K^{t,O}}=\{x:K_{U}^{t,O}(x)\ge |x|-1\}$. Hirahara's prefix-suffix search proof relativizes with exactly the same $O$ throughout: for any hypothetical infinite $P^O$ subset with decider $M^O$, use a retained $n-c\log n$ prefix and enumerate all $n^c$ suffixes. For fixed sufficiently large $c$, the description is shorter than $n-1$ and the runtime is polynomial, eventually below any $t(n)=n^{\omega(1)}$. Hence $R_{K^{t,O}}$ is $P^O$-immune. A hypothetical uniform polynomial-time $V^O(1^n,p,x,\pi)$, complete for every length-$n$ output of $U^O(p)$ with $|p|\le n-2$ in $t(n)$ steps and accepting globally at most two distinct outputs per $(n,p)$ across \emph{all} proofs, yields a language $E^O\in NP^O$ with
\[
|E^O\cap\{0,1\}^n|\le2\sum_{i=0}^{n-2}2^i=2^n-2,\quad
\overline{E^O}\subseteq R_{K^{t,O}}.
\]
For a collapsing oracle $O$ with $P^O=NP^O$, this supplies an infinite $P^O$ subset of a $P^O$-immune language, contradiction. This is a CONDITIONAL no-go for the \emph{universal} schema, not a derivation of ordinary PFCERT or nonalgebrization. In particular do not existentially union exponentially many sparse descriptors, and do not confuse line-query access with explicit expansion of an exponentially large writable configuration.

\paragraph{Full-model and barrier audit.}
All source reading, parsing, gate comparisons, matrices, memory, copying, exponent indexing, computation and output are charged on one execution. The physical state has explicitly $n\le N$ bits, while the numerical horizon can be exponentially large. This does \emph{not} handle a live state with $2^n$ writable bits. DISJ/static rank gate: not a lower-bound argument. Relativization gate: restricted algorithm relativizes, while the universal two-output schema fails matched-oracle collapse. Simultaneous DTISP: positive polynomial time and peak space established, no forbidden SAT point. Configuration counting: no lower bound from the $2^n$ possible states. Algebrization: OPEN. Natural proofs: no large useful circuit property asserted. DEAD/COPY/DELAY/INTERP/RECOMP, affine fast forward and compact circuits are explicit competing controls. INTERP outside the accepted grammar is NOT_CERTIFIED, not hard. Completion-time WRAPPER freshness and short hashes fail to establish classical SAT hardness after disclosure.

\section{TESTS AND OUTCOME}
Five independent local Python checkpoints (seeds 145107,145208,145309,145410,145511) recorded respectively (cases,assertions) $(360,1440),(10,220),(1300,2070),(48,99),(1035,5180)$: exactly \textbf{2,753 cases, 9,009 exact finite assertions, zero final failures}. Code SHA-256: \texttt{e88df3838f64ab1fdaf4ffaa6082f34b7f439452b1a6c657136a2c064a5a3792}. Controls included direct simulation, Boolean M\"obius ANF, malformed or forged certificates, singular affine cores, horizons above $2^{500}$, source-authenticated padded gatewords, SAT/UNSAT reduction and post-disclosure 4/6/8-bit hash-prefix searches. All five stages completed under 30 seconds; these are COMPUTATIONAL rather than general proofs. A six-page full report, companion .bib, exact executable script, raw checkpoint JSON, proposal hash and SHA-256 manifest were preserved locally. The full manuscript compiled without overfull or unresolved reference warnings. Repository checkout preflight was unavailable; GitHub and Drive originals were read through their connected sources. The historical root \texttt{main.tex} remains an explicitly retired placeholder and was not edited.

\paragraph{Five-track status and next target.}
TMR-1: independent pinned counting proof replay and NP-level transfer OPEN. TMR-2: unrestricted actual-source dynamic necessary-work for five-symbol CSP OPEN. TMR-3: arbitrary-machine checkpoint provider and evaluator with complete resource charges OPEN. TMR-4: unrestricted shared/adaptive compiler transfer OPEN. TMR-5: complete independently verified separation OPEN. LSPV-5: universally complete source-grounded PFCERT and nonalgebrization OPEN. \textbf{Next target:} source-discovered, polynomially verifiable noninvertible nonlinear semiconjugacy $P\circ F=G\circ P$ with adaptive writes, using only source-local identities rather than a hidden coNP-hard semantic equality test.

\begin{thebibliography}{9}
\bibitem{hirahara} S. Hirahara. Unexpected Hardness Results for Kolmogorov Complexity Under Uniform Reductions. ECCC TR20-050, 2020. \url{https://eccc.weizmann.ac.il/report/2020/050/}.
\bibitem{bgs} T. Baker, J. Gill, R. Solovay. Relativizations of the P =? NP Question. SIAM J. Comput. 4(4), 1975. DOI: 10.1137/0204037.
\bibitem{bennett} C. H. Bennett. Logical Reversibility of Computation. IBM JRD 17(6), 1973. DOI: 10.1147/rd.176.0525.
\bibitem{toffoli} T. Toffoli. Reversible Computing. MIT-LCS-TM-151, 1980.
\bibitem{aw} S. Aaronson, A. Wigderson. Algebrization: A New Barrier in Complexity Theory. ACM TOCT 1(1), 2009.
\end{thebibliography}
\end{document}
% FINAL_SENTINEL::LSPV5::FACTOREDGATE1::2026-11-18_0706::CONDENSED_PUBLISHED
